% ============================================================================
\section{Algoritmo Genético}\label{sec:genetico}
% ============================================================================

\subsection{Idea}

Los algoritmos genéticos (Holland, 1975) imitan la evolución de las especies. En lugar de una
sola solución que se va modificando, se mantiene una \textbf{población} de soluciones que
evoluciona generación tras generación. El vocabulario es el de la biología:

\begin{center}
\begin{tabular}{@{}lp{11cm}@{}}
\toprule
Término & Significado en optimización \\ \midrule
Individuo & Una solución factible. \\
Gen & Cada una de las componentes con que se codifica la solución. \\
Población & El conjunto de soluciones que se mantiene en cada momento. \\
Aptitud (\emph{fitness}) & Medida de la calidad de un individuo; se deriva del coste. \\
Cruce & Operador que combina dos individuos (padres) para dar otros nuevos (hijos) que
        heredan características de ambos. \\
Mutación & Operador que altera al azar alguna característica de un individuo. \\
Selección & Operador que decide qué individuos sobreviven. \\
Generación & Una iteración del algoritmo. \\
\bottomrule
\end{tabular}
\end{center}

La apuesta del método es que \emph{combinar} soluciones buenas da, con cierta frecuencia,
soluciones aún mejores: si un padre ha resuelto bien una parte del problema y el otro otra
parte, un hijo puede heredar las dos. La mutación introduce material nuevo que el cruce por
sí solo nunca generaría, y la selección empuja a la población hacia las soluciones de menor
coste. Una introducción general gratuita es~\cite[cap.~3]{luke}.

\subsection{Esquema de los apuntes}

\begin{apuntes}[Esquema de los apuntes: Algoritmos Genéticos, algoritmo general]
\begin{tabbing}
\quad\=\kill
Construir una población inicial de soluciones\\
Guardar la mejor solución $x^\ast$\\
Repetir \{\\
\> Seleccionar progenitores y aplicar el OPERADOR DE CRUCE\\
\> Aplicar el OPERADOR DE MUTACIÓN sobre las nuevas soluciones\\
\> Actualizar la mejor solución $x^\ast$ si se ha obtenido alguna mejor\\
\> Aplicar el OPERADOR DE SELECCIÓN\\
\} Mientras no se haya alcanzado el criterio de parada
\end{tabbing}
\end{apuntes}

Los apuntes precisan además cómo entienden cada operador: la población inicial se genera
``por algún algoritmo constructivo''; al aplicar la selección ``generalmente se considera que
mueren la mitad de los individuos, siendo la probabilidad de supervivencia de un individuo
proporcional a su función de aptitud''; los individuos ``se emparejan y dan lugar a varios
hijos''; y la mutación ``se aplica con una probabilidad muy baja''. El diseño siguiente
respeta esas indicaciones.

\subsection{Adaptación al problema de cubrimiento}

\paragraph{Representación.} Los apuntes subrayan que representar bien las soluciones es la
decisión principal.

\begin{decision}[Decisión: representación binaria, guardada como conjunto]
Un individuo es el vector binario $x$ del modelo: un gen por columna, que vale 1 si la
columna está elegida. Es la representación natural del problema y hace que cruzar y mutar
sean operaciones evidentes. Como de los 10\,000 genes solo unos 150 valen 1, en memoria no
se guarda el vector entero sino el conjunto de columnas elegidas, con la misma estructura
\cod{Solucion} que usan los demás algoritmos.
\end{decision}

\paragraph{Aptitud.} Se minimiza el coste, así que es más apto quien menos cuesta. Para un
individuo de coste $z$ en una población cuyos costes van de $z_{\min}$ a $z_{\max}$:
\[
   \text{aptitud} \;=\; (z_{\max}-z) \;+\; 0.05\,(z_{\max}-z_{\min}) \;+\; 10^{-6}.
\]
El primer término es la ventaja sobre el peor; el segundo es un pequeño margen para que el
peor individuo no tenga aptitud cero (y conserve alguna posibilidad de sobrevivir). Responde
a la observación de los apuntes~\cite{yanez}: ``en el caso de problemas de minimización, hay
que modificar la función de adaptación para que sean más probables los de menor valor''.

\paragraph{Población inicial.} El primer individuo es la solución del voraz puro; los demás
se construyen con el voraz aleatorizado (LRC de tamaño 5) para que sean distintos entre sí.
Responde a la indicación de los apuntes sobre ``cómo construir conjuntos de soluciones
diversas'': si todos los individuos fueran iguales, el cruce no produciría nada nuevo.

\paragraph{Progenitores.} Se baraja la población y se empareja de dos en dos. Cada pareja
tiene dos hijos, de modo que en cada generación nacen tantos hijos como individuos hay.

\begin{decision}[Decisión: todos se reproducen; la presión se ejerce al sobrevivir]
No se favorece a los mejores a la hora de elegir progenitores porque la presión selectiva ya
se aplica en el operador de selección (que, según los apuntes, es el que elimina individuos).
Aplicarla dos veces haría que la población se volviera uniforme demasiado pronto
(\emph{convergencia prematura}).
\end{decision}

\paragraph{Cruce.} Los genes en los que ambos padres coinciden se heredan tal cual: si los
dos tienen la columna, los dos hijos la tienen; si ninguno la tiene, ninguno. Los genes en
que difieren se reparten, y los dos hijos son complementarios (lo que no hereda uno de un
padre lo hereda el otro). Hay tres formas de repartir (parámetro \cod{cruce}):
\begin{enumerate}
  \item \textbf{Un punto}: se sortea una columna de corte; el primer hijo toma del padre los
        genes anteriores al corte y de la madre los posteriores, y el segundo al revés. Es el
        cruce del ejemplo de los apuntes~\cite{yanez} (``uno con los 3 primeros bits del padre
        y los 2 últimos de la madre y el otro con los 3 primeros bits de la madre y los otros
        2 del padre'').
  \item \textbf{Uniforme}: cada gen se toma de uno u otro progenitor a cara o cruz.
  \item \textbf{Ponderado por coste (``fusión'')}: como el uniforme, pero el gen se toma del
        progenitor \emph{más barato} con mayor probabilidad: si los costes son $z_a$ y
        $z_b$, se hereda de $a$ con probabilidad $z_b/(z_a+z_b)$.
\end{enumerate}
La tercera variante se inspira en el operador de fusión que Beasley y Chu propusieron para
este mismo problema, mencionado en~\cite{tayebi}\footnote{El artículo original (J.~E.~Beasley
y P.~C.~Chu, \emph{A genetic algorithm for the set covering problem}, European Journal of
Operational Research 94, 1996) es de pago y no se ha utilizado; aquí solo se toma la idea
general de ponderar la herencia por la aptitud, que se describe en fuentes de acceso libre.}.

\paragraph{Reparación.} Un hijo puede dejar filas sin cubrir.

\begin{decision}[Decisión: reparar a los hijos]
Tras el cruce y la mutación, cada hijo se completa con el voraz (\cod{completar}) y se le
quitan las columnas redundantes. Así todos los individuos son siempre soluciones factibles y
sin columnas sobrantes, como en los demás algoritmos. La alternativa (admitir individuos
infactibles y penalizarlos en la aptitud) obligaría a calibrar la penalización y llenaría la
población de individuos inútiles.
\end{decision}

\paragraph{Mutación.} Con probabilidad \cod{pmut}, un hijo sufre una mutación que altera
\cod{nmut} genes.

\begin{decision}[Decisión: mutación equilibrada entre encender y apagar]
Si el gen a alterar se eligiera al azar entre los 10\,000, casi siempre sería un 0 (solo el
1.5\,\% valen 1): la mutación casi siempre añadiría una columna, que además sería redundante
y la reparación eliminaría enseguida. No cambiaría nada. Por eso cada alteración es, con
probabilidad $1/2$, apagar un gen encendido (quitar una columna de la solución, lo que
obliga a la reparación a buscar otra forma de cubrir sus filas) y, con probabilidad $1/2$,
encender un gen apagado al azar.
\end{decision}

\paragraph{Mejora local (opcional).} Con \cod{bl} $>0$ se aplica a cada hijo el VND con
\cod{bl} entornos. Un genético con búsqueda local se suele llamar \emph{algoritmo memético}
(figura como tal en la clasificación de los apuntes). Es opcional para poder medir en la
calibración cuánto aporta.

\paragraph{Selección.} Al final de cada generación hay $2\cdot\cod{pob}$ individuos (padres e
hijos) y sobrevive la mitad, como indican los apuntes:
\begin{itemize}
  \item el mejor sobrevive siempre (\textbf{elitismo}): la mejor solución no se pierde nunca;
  \item el resto de plazas se sortean una a una con probabilidad proporcional a la aptitud
        (método de la \textbf{ruleta}). A diferencia del ejemplo de los apuntes, que permite
        la repetición, aquí un individuo no puede ser elegido dos veces: con poblaciones
        pequeñas, permitir repeticiones llena enseguida la población de copias del mejor;
  \item los individuos \emph{repetidos} (copias exactas de otro) reciben aptitud cero, para
        que la población no se llene de clones.
\end{itemize}

\subsection{Parámetros}

\begin{center}
\begin{tabular}{@{}lp{11.3cm}@{}}
\toprule
Nombre & Significado y efecto \\ \midrule
\cod{pob} & Tamaño de la población. Grande: más diversidad, pero menos generaciones en el
            mismo tiempo. \\
\cod{cruce} & Tipo de cruce: 1 = un punto, 2 = uniforme, 3 = ponderado por coste. \\
\cod{pmut} & Probabilidad de que un hijo mute. Con 0 no hay mutación. \\
\cod{nmut} & Número de genes que altera una mutación. \\
\cod{bl} & Entornos de la mejora local de los hijos (0 = sin mejora local). \\
\bottomrule
\end{tabular}
\end{center}

\subsection{Recorrido por el código}

El cruce recorre los genes que valen 1 en cada padre (los que valen 0 en ambos no hay que
tocarlos). La variable \cod{de\_a} decide de quién hereda el primer hijo:

\fragmento{genetico.c}{64}{85}

La mutación:

\fragmento{genetico.c}{95}{107}

El cuerpo de una generación (cruce, mutación, reparación, mejora local y actualización de
$x^\ast$) sigue el orden del esquema:

\fragmento{genetico.c}{165}{183}

Y la selección por ruleta con elitismo. El primer superviviente ($t=0$) es el mejor; los
demás se sortean entre los que aún no han sido elegidos:

\fragmento{genetico.c}{204}{225}
